%==============================================================================
% Section 9: Discussion and outlook
%==============================================================================
\section{Discussion and outlook}
\label{sec:discussion}

\subsection{What has been established, by layer}
\label{sec:discussion:layers}

The results of this paper can be organised along their logical layers of
establishment.  All of the following are claims under the DBR protocol.

\paragraph{Quadratic level (Sections~\ref{sec:rankdefect} and~\ref{sec:harmonic}).}
In the DBR quadratic constrained spectrum, we proved that the dimension of the
non-propagating scalar kernel is realised as
$\dzero=\dim\ker n=3-\rank(n)$.  We also showed that two TT modes propagate with
positive kinetic, that $\dzero$ is invariant under SO(3) rotations of the spatial
frame, and the activation-rank mechanism.  At this layer we do not identify the
kernel with propagating zero modes, do not claim $\dzero$ as a new invariant, and
do not extend to class B.

\paragraph{$k>0$ spectral completion (Section~\ref{sec:spectral}).}
In the analysed harmonic sectors we confirmed the propagation of the 2 TT, the
no-ghost property, and the non-contamination of $\dzero$, through the endpoint
controls ($\Tthree,\Sthree$), the closed-form completion of $\Nilthree$, and the
bulk gap and positive compact edge of $\Solthree$.  This is a robustness
restricted to the analysed harmonic sectors (the $\Solthree$ edge is
representative numerics), and is not a claim of full PDE closure, an all-Anosov
closed-form edge theorem, or a compact continuous spectrum.

\paragraph{Classification of finite deformations (Section~\ref{sec:finite}).}
Within the finite Palatini branch / homogeneous minisuperspace we exhibited the
modulus/wall split ($\Tthree$ 3/0, $\Nilthree$ 1/1, $\Solthree$ 0/1, $\Sthree$
0/0) and the minisuperspace no-ghost.  This is a supporting result and does not
include a nonlinear lift/mass theorem or an identification of the moduli count
with $\dzero$.

\paragraph{Torsion closure (Section~\ref{sec:torsion}).}
Within the inspected $\AX/\VT/\MX$ ansatz, we showed that the reduced torsion
sector remains an exact auxiliary shell, generates no quadratic kinetic term
under field promotion, and leaves the finite-deformation classification
invariant.  This is not a full 24-component $(\partial T)^{2}$-absence theorem, a
general claim of NY bulk non-involvement, or a charge/localisation theorem.

\medskip
Separately, the route comparison (Section~\ref{sec:dbr},
Appendix~\ref{app:route}) shows, as a control, that the RBD route produces a
Dirac-inequivalent extra mode in the $\Sthree$ scalar sector.  This is the basis
for adopting DBR and does not claim a general failure of RBD or a physical-ghost
interpretation.

\subsection{Scope of the novelty (restated)}
\label{sec:discussion:novelty}

The topology hierarchy $\Sthree{:}0$ / $\Solthree{:}1$ / $\Nilthree{:}2$ /
$\Tthree{:}3$ of this paper is, algebraically, the \emph{known data} of the
rank-defect / abelianisation / $b_{1}$ of the Bianchi class-A structure
matrix~\cite{ref16,ref41,ref42}.  Our contribution is not its discovery but the three
points: (i) the verification that, in the Lorentzian first-order EC+NY system,
the DBR constraint dynamics \emph{independently reproduces} this data as a
residual non-propagating scalar kernel; (ii) the four-topology comparison under a
single protocol; and (iii) that this classification is protected by the
three-way robustness checks of the $k>0$ spectral completion, the
finite-deformation classification, and the torsion closure.

\subsection{Relation to Teichm\"uller/moduli counts of compact Bianchi models}
\label{sec:discussion:teichmuller}

The global degrees of freedom (Teichm\"uller parameters) of compact Bianchi
models have been systematically counted by
Koike--Tanimoto--Hosoya~\cite{ref24,ref25,ref26}.  Our $\dzero$ is a
\emph{different object}: the KTH moduli are global structural deformations of the
compact quotient of metric GR, whereas our kernel is the non-propagating
direction of the homogeneous perturbations in the DBR constrained spectrum of the
first-order EC+NY system.  Indeed, internally to this paper, the finite exact
moduli count ($\Tthree$:3, $\Nilthree$:1, $\Solthree$:0) does not in general equal
$\dzero$ (3, 2, 1) (Section~\ref{sec:finite}).  The precise correspondence among
the quadratic kernel, the finite moduli, and the Teichm\"uller deformations
(which kernel directions lift to global deformations) remains a task for future
work.

\subsection{Outlook}
\label{sec:discussion:outlook}

\begin{itemize}
  \item \textbf{Higher-order lift}: a full nonlinear lift/mass analysis beyond
        the finite Palatini branch (to settle the dynamical meaning of the wall).
  \item \textbf{Extension to class B / VII$_{0}$ / VIII}: our formula is
        restricted to class A.  In class B there are known obstructions to the
        symmetry reduction itself~\cite{ref21}, requiring a separate treatment.
        It is also known that the Mathieu structure appears in the spectral
        problem of the Bianchi VII system.
  \item \textbf{Closed form for the $\Solthree$ edge}: an upgrade of the compact
        edge $\lambda_{1}$, kept at representative numerics, to an all-Anosov
        closed-form theorem.
  \item \textbf{Correspondence with the Euclidean series}: a comparison of the
        results of this paper (Theorem~\ref{thm:rankdefect},
        Table~\ref{tab:instances}, Table~\ref{tab:modulus},
        Table~\ref{tab:torsion}) with the Euclidean response
        dictionary~\cite{ref4,ref5}.
\end{itemize}
